% A.121: full octic common-field classification; root owns integration.
\subsection{The complete common abelian field of an eighth-root fibre}

\begin{statement}{A.121}{Proved}
{All common octic characters from explicit fourth-power relations}
Let $D_1,\ldots,D_k\in\Q^\times$ and put
\[
 \mathcal E_8=\Q[Z_1^{\pm1},\ldots,Z_k^{\pm1}]
                      /(Z_i^8-D_i)_i,\qquad
 K=\Q(\sqrt{D_1},\ldots,\sqrt{D_k}).
\tag{A.121.1}
\]
Let $B_4$ be the exact common abelian field of the quartic
fibre computed in A.96.
Write
\[
 D_i=(-1)^{\delta_i}\prod_{P\in\mathcal P}P^{w_{iP}},
 \qquad \delta_i\in\{0,1\},\qquad
 F=\Q(\mu_{16},\sqrt P:P\in\mathcal P).
\tag{A.121.2}
\]
The largest abelian field $B_8$ contained in every field
factor of $\mathcal E_8$ is given by the following exhaustive
three cases:
\[
\begin{array}{ll}
 i\notin K:& B_8=K=B_4,\\
 i\in K,\ \sqrt2\notin B_4:& B_8=B_4,\\
 \mu_8\subset B_4:&
 B_8=B_4\bigl(K_t:t\in\mathcal R_4\bigr).
\end{array}
\tag{A.121.3}
\]
Here the finite relation module and the actual rational
numbers defining its character fields are
\[
 \begin{gathered}
 \mathcal R_4=
 \left\{t\in(\Z/4)^k:
            \sum_i t_iw_{iP}=0\pmod4\ \text{for every }P\right\},\\
 t_i\in\{0,1,2,3\},\qquad
 \delta_t=\sum_i t_i\delta_i\pmod2,\qquad
 r_t=\prod_P P^{\,\frac14\sum_i t_iw_{iP}}>0,\\
 \prod_iD_i^{t_i}=(-1)^{\delta_t}r_t^4 .
 \end{gathered}
\tag{A.121.4}
\]
For $\delta_t=0$, set $K_t=\Q(\sqrt{r_t})$.
For $\delta_t=1$, let $K_t$ be the cyclic quartic field
cut out by
\[
             \Psi_{r_t}=\chi_{16}+\chi_{r_t},
 \qquad
 \chi_{16}(5)=1/4,\quad\chi_{16}(-1)=0
                         \quad\text{on }(\Z/16)^\times.
\tag{A.121.5}
\]
In particular (A.121.3)--(A.121.5) compute every common
character, not only a collection of characters contained
in the common field.  The inclusion of each displayed
generator into every actual field factor is proved below.
In the third case $B_8/B_4$ is an elementary abelian
two-extension; the characters over $\Q$ can have order four.

For a concrete harmonic fibre, take
\[
 R=\frac17,\quad c=-2,\quad
 (D_0,D_1,D_{1/2})=(-14,7/2,7).
\tag{A.121.6}
\]
Its full common field is
$B_8=\Q(\mu_{16},\sqrt7)$, of degree $16$.
The mixed relation $D_0D_1D_{1/2}^{-2}=-1$ gives
the pure character $\chi_{16}$ by the actual root
quotient $(Z_0Z_1)/Z_{1/2}^2$.
Thus the negative-fourth-power triple relation of
A.116 is one specified family within the complete
list (A.121.4).
The field calculation alone makes no assertion that
all characters in this complete list have constant
evaluation on the corresponding old eighth-power cover.
\end{statement}
\begin{proof}
\emph{The actual finite group and its root stabilizers.}
Use the fields and restriction map of A.57 at $n=8$:
\[
 L=F(P^{1/8}:P\in\mathcal P),\qquad
 G=\operatorname{Gal}(L/\Q),\qquad
 \pi:G\twoheadrightarrow\Gamma=\operatorname{Gal}(F/\Q).
\tag{A.121.7}
\]
The positive real-root degree calculation there gives
\[
 \operatorname{Gal}(L/F)=(\Z/4)^{\mathcal P},
 \qquad [G,G]=\operatorname{Gal}(L/F).
\tag{A.121.8}
\]
For $\sigma\in\Gamma$, retain the actual compatible
cyclotomic and quadratic signs
\[
 \sigma(\zeta_{16})=\zeta_{16}^{u},\qquad
 \sigma(\sqrt P)=(-1)^{\epsilon_P}\sqrt P.
\]
In particular the sign on $\sqrt2$, when it occurs,
is determined by $u$; it is not an extra independent
sign variable.
All lifts of $\sigma$ are given by
\[
 P^{1/8}\longmapsto\zeta_8^{\,k_P}P^{1/8},
       \qquad k_P=\epsilon_P+2h_P,\quad h_P\in\Z/4.
\tag{A.121.9}
\]
The root tuples are
$z_i=\zeta_{16}^{\delta_i+2e_i}\prod_P P^{w_{iP}/8}$,
where $e_i\in\Z/8$.
A lift in (A.121.9) fixes this tuple exactly when
\[
 (u-1)(\delta_i+2e_i)
       +2\sum_Pw_{iP}\epsilon_P
       +4\sum_Pw_{iP}h_P=0\pmod {16}
                                     \quad(1\le i\le k).
\tag{A.121.10}
\]
Let $\mathcal A$ be the union of the images
$\pi(G_\omega)$ of these actual root stabilizers, and put
$H=\langle\mathcal A\rangle$.
Equation (A.121.8) and the finite fixed-field maps give
\[
                         B_8=F^H .
\tag{A.121.11}
\]
This is the same passage from all field factors to
their common abelian field as in A.57.

The square-root maps $z_i\mapsto z_i^4$ give $K\subset B_8$.
The square-coordinate map from the octic fibre to the quartic
fibre, $z_i\mapsto z_i^2$, also gives $B_4\subset B_8$:
each octic field factor contains the actual quartic factor
selected by its squared tuple.
Therefore
\[
              H\subset\Gamma_4:=\operatorname{Gal}(F/B_4)
                          \subset\Gamma_K:=\operatorname{Gal}(F/K).
\tag{A.121.12}
\]

\emph{The case without a square root of minus one in $K$.}
Here is the specialization of A.80's exact argument.
For $\sigma\in\Gamma_K$ with $u\equiv3\pmod4$,
$\gcd(8,u-1)=2$.
Choose any $k_P$ with the parity $\epsilon_P$.
The condition that $\sigma$ fixes every $\sqrt{D_i}$ is
\[
        \frac{u-1}{2}\delta_i+\sum_Pw_{iP}k_P=0\pmod2.
\]
It is exactly the divisibility condition allowing one to
solve (A.121.10) for each $e_i$.
Thus $\sigma$ is the image of an actual root stabilizer.
Since $i\in F$ and $i\notin K$, there is
$\tau\in\Gamma_K$ with $u_\tau\equiv3\pmod4$.
Every element of $\Gamma_K$ either has this same
cyclotomic parity or becomes one after multiplication by
$\tau$.
Both factors then belong to $H$, proving $H=\Gamma_K$.
The fixed-field identity (A.121.11) proves $B_8=K$.
The same hypothesis gives $B_4=K$ by A.96.

\emph{Reduction of the remaining stabilizer equations.}
Assume $i\in K$.
Every element of $\Gamma_K$ then has
$u\in\{1,5,9,13\}$ and
$\sum_Pw_{iP}\epsilon_P$ even.
Define the integral vector
\[
       A_i(\sigma)
       =\frac12\sum_Pw_{iP}\epsilon_P
                    +\frac{u-1}{4}\delta_i .
\tag{A.121.13}
\]
After division by four, (A.121.10) becomes
\[
       \sum_Pw_{iP}h_P+\frac{u-1}{2}e_i
                      =-A_i(\sigma)\pmod4 .
\tag{A.121.14}
\]
Here $W=(w_{iP})$ is the integral valuation matrix.
When $u\equiv5\pmod8$, the coefficient of $e_i$
has image $2\Z/4$.
Consequently (A.121.14) is soluble exactly when
\[
                             Wh=-A(\sigma)\pmod2.
\tag{A.121.15}
\]
This is precisely the quartic stabilizer system
(A.96.7), with $u$ reduced modulo eight.
A.96 proved that its solution set is the full subgroup
fixing $B_4$, rather than merely a set generating that
subgroup.
We therefore have the exact equality
\[
 \mathcal A\cap\{u\equiv5\pmod8\}
          =\Gamma_4\cap\{u\equiv5\pmod8\}.
\tag{A.121.16}
\]
In the forward direction this also follows from the
actual square-coordinate map; in the reverse direction
a solution of (A.121.15) is lifted to $h_P\in\Z/4$,
and each $e_i$ is chosen to correct the remaining even
error in (A.121.14).
This supplies an actual lift and root tuple.

\emph{The middle case and the subgroup it generates.}
Suppose $\sqrt2\notin B_4$.
Since $i\in B_4\subset F$, restriction to $\Q(\sqrt2)$
gives an element $\tau\in\Gamma_4$ acting nontrivially
on $\sqrt2$, and hence satisfying $u_\tau\equiv5\pmod8$.
Equation (A.121.16) gives $\tau\in\mathcal A\subset H$.
For arbitrary $\sigma\in\Gamma_4$, either
$u_\sigma\equiv5\pmod8$, or
$u_{\sigma\tau}\equiv5\pmod8$.
Equation (A.121.16) then puts respectively $\sigma$, or
both $\sigma\tau$ and $\tau$, in $H$.
Thus $\Gamma_4\subset H$; the reverse inclusion is
(A.121.12).
Taking the exact fixed fields proves $B_8=B_4$.
This argument uses the subgroup generated by stabilizer
images; it does not require their union to be a subgroup
in this middle case.

\emph{All remaining characters when $\mu_8\subset B_4$.}
Since $i\in B_4$, the remaining case
$\sqrt2\in B_4$ is exactly $\mu_8\subset B_4$.
Now every $\sigma\in\Gamma_4$ has $u\in\{1,9\}$.
The term involving $e_i$ in (A.121.14) is zero.
Hence the exact condition is
\[
                         Wh=-A(\sigma)\pmod4 .
\tag{A.121.17}
\]
The perfect dot pairing on $(\Z/4)^k$ identifies the
annihilator of the image of $W$ with
$\mathcal R_4=\ker(W^t\bmod4)$.
In particular (A.121.17) is soluble if and only if
\[
                  \sum_i t_iA_i(\sigma)=0\pmod4
                                  \quad(t\in\mathcal R_4).
\tag{A.121.18}
\]
For completeness, this orthogonality equivalence follows
by the elementary row and column operations over $\Z/4$
which put a matrix into a diagonal matrix with entries
$1,2,0$.
A target coordinate for a diagonal $1$ is unrestricted;
for a $2$ it must be even, exactly its pairing condition
against the annihilator $2$; for a $0$ it must vanish,
exactly its pairing condition against the full
$\Z/4$.
The invertible operations transport the dot pairings
and prove (A.121.18) for the original matrix.

For $t\in\mathcal R_4$, write
$v_P=\frac14\sum_i t_iw_{iP}$, so $r_t=\prod_PP^{v_P}$.
Equations (A.121.13) and $u=1,9$ give
\[
 \begin{split}
 \sum_i t_iA_i(\sigma)
    &=2\sum_Pv_P\epsilon_P
                     +\frac{u-1}{4}\delta_t\pmod4\\
    &=4\bigl(\chi_{r_t}(\sigma)
                          +\delta_t\chi_{16}(\sigma)\bigr)
                                                 \pmod4.
 \end{split}
\tag{A.121.19}
\]
Here the last expression means the representative in
$\Z/4$ under $\tfrac14\Z/\Z\simeq\Z/4$.
Indeed the two possible values of the integral
cyclotomic character $4\chi_{16}(u)$ are $0,2$,
which equal $(u-1)/4$ for $u=1,9$.
Thus every equation in (A.121.18) is the vanishing
of the literal restriction of the character defining $K_t$.

It follows that in this last case the union of the
stabilizer images is already the subgroup
\[
 \mathcal A=
 \Gamma_4\cap
    \bigcap_{t\in\mathcal R_4}
       \ker\bigl(\chi_{r_t}+\delta_t\chi_{16}\bigr).
\tag{A.121.20}
\]
The reverse inclusion supplies actual $h_P$ through
(A.121.17) and arbitrary $e_i$ in (A.121.14);
there is therefore an actual root tuple fixed by the
resulting lift.
Equations (A.121.11) and (A.121.20) prove the third
line of (A.121.3) and exhaust every common character.
All the characters in (A.121.20) restrict to characters
of order at most two on $\Gamma_4$, proving the final
assertion about $B_8/B_4$.

\emph{The displayed fields inside every individual factor.}
The fixed-field calculation also has the following
direct root maps.
For $t\in\mathcal R_4$, take the actual product
$z_t=\prod_i z_i^{t_i}$ inside any one octic field
factor.
It satisfies
\[
                         z_t^8=(-1)^{\delta_t}r_t^4.
\tag{A.121.21}
\]
If $\delta_t=0$, it is of the form
$\zeta_8^j\sqrt{r_t}$.
Every factor already contains $\mu_8\subset B_4$,
so division by that root of unity puts
$\sqrt{r_t}$ in this same factor.
If $\delta_t=1$, it is of the form
$\zeta_{16}^j\sqrt{r_t}$, with $j$ odd.
An automorphism fixing this root has
$u=1+8\epsilon\pmod{16}$, where $\epsilon$
is its sign on $\sqrt{r_t}$.
Consequently
$\chi_{16}+\chi_{r_t}=\epsilon/2+\epsilon/2=0$.
This proves that the field $K_t$ is contained in the
actual field generated by $z_t$, exactly as in
(A.116.3).
These are explicit inclusions into each field factor,
compatible with (A.121.20).

\emph{The harmonic example and a positive fourth-power direction.}
For (A.121.6), the individual quadratic field is
\[
 K=\Q(\sqrt{-14},\sqrt{7/2},\sqrt7)
       =\Q(i,\sqrt2,\sqrt7).
\tag{A.121.22}
\]
Indeed the last two radicals give $\sqrt2$, and
their combination with the first gives $i$.
The only nontrivial negative-square subset is
$D_0D_1=-49$; its quartic extra generator is
$\sqrt{14}$, already in $K$.
Thus $B_4=K$.
The root quotient in the statement has eighth power
$-1$, so it is a primitive sixteenth root of unity.
Every factor therefore contains
$\Q(\mu_{16},\sqrt7)=F$.
The upper bound (A.121.11) gives equality, and the
degree is $8\cdot2=16$ because $\sqrt7\notin\Q(\mu_{16})$.
For example the latter follows from ramification:
$7$ ramifies in $\Q(\sqrt7)$ and not in
$\Q(\mu_{16})$.

There can also be additional quadratic characters
arising from positive fourth-power relations.
For
\[
                   (D_1,D_2,D_3)=(-9,625,2)
\tag{A.121.23}
\]
one has $B_4=\Q(i,\sqrt2,\sqrt3)$.
The relation $D_2=5^4$ supplies $\sqrt5$ through
(A.121.21).
Every negative fourth-power relation would have an
odd exponent of $D_1$, whereas its valuation at
three would then be $2$ modulo four; hence no such
relation exists.
The positive relations in (A.121.4) add only
$\sqrt5$ modulo $B_4$.
Thus
\[
                  B_8=\Q(i,\sqrt2,\sqrt3,\sqrt5),
                              \qquad [B_8:\Q]=16.
\tag{A.121.24}
\]
This second example verifies why positive relations
must be retained in the complete formula.
\end{proof}

The exact certificate
\path{code/eighth_common_field_certificate.py}
compares (A.121.3) with the literal root-stabilizer
system (A.121.10), retaining the cyclotomic sign
of $\sqrt2$ and both subgroup-generation steps.
Its output is saved beside the program.
The separate search
\path{code/harmonic_eighth_remaining_search.py}
retains twenty-six harmonic fibres with extra character
directions beyond the earlier specified norm families.
It records 39074 trials and 6136 actual local eighth-power
points without a nonzero evaluation in its sampled range.
Those finite evaluations do not prove constancy for
the full character list (A.121.4).

